\documentclass[11pt,a4paper]{article}
\usepackage[utf8]{inputenc}
\usepackage[T1]{fontenc}
\usepackage{lmodern}
\usepackage[margin=24mm]{geometry}
\usepackage{amsmath,amssymb,mathtools}
\usepackage{booktabs}
\usepackage{enumitem}
\usepackage{xcolor}
\usepackage[hidelinks]{hyperref}

\newcommand{\dd}{\mathop{}\!\mathrm{d}}
\newcommand{\ii}{\mathrm{i}}
\newcommand{\id}{\mathrm{id}}
\newcommand{\R}{\mathbb{R}}
\newcommand{\Z}{\mathbb{Z}}
\newcommand{\N}{\mathbb{N}}
\newcommand{\Ord}{\operatorname{ord}}
\newcommand{\Ker}{\operatorname{Ker}}
\newcommand{\Aut}{\operatorname{Aut}}
\newcommand{\slashed}[1]{\not\!#1}
\setlength{\parindent}{0pt}
\setlength{\parskip}{6pt}
\allowdisplaybreaks

\title{Particle Physics --- Decay Rates}
\author{IMAPP lecture notes}
\date{Spring 2026}
\begin{document}\maketitle

\section{Exponential decay}
If $N(t)$ particles are present, a constant decay rate gives
\[
\dd N=-\Gamma N\,\dd t,\qquad N(t)=N(0)e^{-\Gamma t},\qquad \tau=\Gamma^{-1}.
\]
The width $\Gamma$ is the natural width associated with the finite lifetime.

\section{Decay width}
For a decay into a set of final states, the differential width is written schematically as
\[
\dd\Gamma=\frac{1}{2E_A}|\mathcal M|^2\dd\Phi_n,
\]
where $\mathcal M$ is the Lorentz-invariant matrix element and $\dd\Phi_n$ is the Lorentz-invariant $n$-body phase-space measure. For two final particles,
\[
\dd\Phi_2=(2\pi)^4\delta^{(4)}\!\left(p_A-\sum_i p_i\right)
\prod_{i=1}^2\frac{\dd^3\mathbf p_i}{(2\pi)^3 2E_i}.
\]
The total width is obtained by integrating over all allowed final-state momenta.

\section{Cross-sections}
The interaction rate is the incident flux multiplied by the number of target particles and the cross-section:
\[
\text{rate}=\text{flux}\times N_{\text{target}}\times\sigma.
\]
The cross-section has dimensions of area; the barn is a standard unit, $1\,\mathrm{b}=10^{-28}\,\mathrm{m}^2$.

For two incoming particles, the Lorentz-invariant flux can be expressed in terms of their four-momenta. The invariant matrix element and phase-space factors together ensure that the calculated rate is frame independent.

\section*{Summary}
Decay widths describe the probability per unit time for an unstable state to decay. Cross-sections describe interaction probabilities per incident flux. Both are built from invariant amplitudes and Lorentz-invariant phase space.
\end{document}
